\documentclass{article}

\usepackage{amsmath,amssymb,amsthm,mathtools,bm}
\usepackage{booktabs}
\usepackage{graphicx}
\usepackage{microtype}
\usepackage{hyperref}
\usepackage{icml2026}

\newtheorem{theorem}{Theorem}
\newtheorem{proposition}[theorem]{Proposition}
\newtheorem{lemma}[theorem]{Lemma}
\newtheorem{corollary}[theorem]{Corollary}
\newcommand{\R}{\mathbb R}
\newcommand{\E}{\mathbb E}
\newcommand{\norm}[1]{\left\lVert #1\right\rVert}
\newcommand{\Proj}{\operatorname{Proj}}
\newcommand{\dist}{\operatorname{dist}}

\icmltitlerunning{Constraint-Aware Multilevel Picard Methods}

\begin{document}

\twocolumn[
\icmltitle{Constraint-Aware Multilevel Picard Methods}

\begin{icmlauthorlist}
\icmlauthor{Yeping Wang}{gatech}
\end{icmlauthorlist}
\icmlaffiliation{gatech}{H. Milton Stewart School of Industrial and Systems Engineering, Georgia Institute of Technology, USA}
\icmlcorrespondingauthor{Yeping Wang}{anonymous@anonymous.edu}
\icmlkeywords{multilevel Picard, high-dimensional PDE, Monte Carlo, hard constraints, invariant regions}
\vskip 0.3in
]
\printAffiliationsAndNotice{}

\begin{abstract}
Multilevel Picard (MLP) methods repeatedly feed finite-sample value--gradient estimates into a nonlinear generator. We study a simple failure mode: a noisy child state can leave a region that PDE structure certifies for the exact solution, and the generator can amplify that infeasible component before Monte Carlo averaging has a chance to cancel it. We propose constraint-aware MLP, which retracts child states into certified feasible sets immediately before nonlinear reuse, with both samplewise nearest-point and batch-coupled gauge realizations. On a known high-dimensional MLP counterexample, a certified active subspace makes the offending generator vanish; recursively enforcing that subspace removes every nonlinear correction pathwise and yields exact full-state MSE $(4-2/\pi)/N$, independent of ambient dimension. Broader HJB, gradient-dependent, and nonlinear-finance experiments show substantial finite-budget gains and a useful samplewise--batchwise bias--variance tradeoff, while negative controls show little effect when the constraint mechanism is inactive.
\end{abstract}

\section{Introduction}
High-dimensional MLP methods combine Picard iteration with Monte Carlo and multilevel differences to solve semilinear parabolic PDEs without spatial grids \citep{E2021,HutzenthalerKruse2020,NeufeldWu2025}. Their recursive structure creates a numerical issue that is easy to miss* a finite-sample value--gradient state is itself an input to later nonlinear evaluations. Thus an error that would be harmless under linear averaging can become large after a norm, square, threshold, or other nonlinear transformation.

Suppose PDE structure certifies that the exact state $Y^\star=(u,z)$ lies in a set $\mathcal C(t,x)$, but a Monte Carlo child $\widehat Y$ lies outside. The usual recursion evaluates $f(t,x,\widehat Y)$ anyway. Our proposal is minimal:
\begin{equation}
\widehat Y \;\longrightarrow\; \widetilde Y=\mathcal R(\widehat Y)\in\mathcal C(t,x)
\;\longrightarrow\; f(t,x,\widetilde Y),
\label{eq:principle}
\end{equation}
where the correction occurs \emph{inside} the tree, before nonlinear reuse. We call this constraint-aware MLP.

We study two realizations. \textbf{Samplewise IR-MLP} applies the metric projection independently to each child. \textbf{Batch-IR} uses one common gauge factor for a sibling Monte Carlo batch. Samplewise projection has cleaner geometry; batchwise contraction couples samples and moves the empirical mean, but can suppress batch spread more aggressively.

The paper is intentionally not a new general convergence theory for MLP. We use only enough standard theory to establish that the samplewise correction preserves the intended target under the usual assumptions; a concise inheritance argument is placed in Appendix~\ref{app:convergence}. The main theoretical contribution is instead an exact structural rescue of the dimensionality counterexample of \citet{HutzenthalerNguyen2025}. There the true gradient lies in $\operatorname{span}(e_1)$, while the driver depends only on the other coordinates. Recursively returning to that certified subspace makes every nonlinear correction vanish pathwise and gives a dimension-free error formula.

\paragraph{Contributions.}
(1) We formulate recursive feasible-state correction for MLP and compare samplewise and batch-coupled retractions. (2) We prove an exact generator-compatible rescue theorem and dimension-free MSE on a published counterexample. (3) We connect the geometric mechanism to high-dimensional HJB, Neufeld--Wu, and nonlinear-funding experiments, with negative controls. (4) We keep the generic convergence claim deliberately narrow: target preservation and standard samplewise inheritance are defensive results; a sharp theory for batch-coupled retractions is left open.

